\section{Proofs and supplements for Section~\ref{sec:limits}}\label{app:limits}

This appendix computes the conditioning of the bounded-coefficient reduction
of Del Pia and Khajavirad, proves Proposition~\ref{prop:lbproduct} and the
propositions of Sections~\ref{sec:limits-messages}--\ref{sec:limits-constraints},
and gives the local moment relaxations of Section~\ref{sec:limits-local}.

\subsection{Conditioning of the bounded-coefficient reduction}\label{app:dk}

\begin{remark}[Conditioning of the bounded-coefficient reduction]
\label{lim:rem:dk}
Consider the construction of \cite[Theorem~3]{DelPiaKhajavirad2026} for the
two-item instance $a_1=c$, $a_2=c-1$ with target $a_0=c$, where $c\ge3$, with
objective $\Psi$ (their (20)--(24), all weights one) on the unit box. In
their notation, $\ell=\lceil\log_2(2c)\rceil$, $D=2^\ell\ge2c$, and there are
$5\ell+2$ variables. Every term of $\Psi$ is nonnegative. A zero of $\Psi$
encodes a binary solution, and $(1,0)$ is the only one; all copied and state
variables are then determined by the recurrences. Hence $\Psi$ has a unique
minimizer $v^*$ with value zero, and by Lemma~\ref{lem:unique-growth} some
growth constant $g>0$ is valid.

Encode instead the choice $(0,1)$ with every copy, digit and running-sum
recurrence exact. Only the final square
$(s_2-w_\ell)^2=((c-1-c)/D)^2$ is nonzero, so $\Psi(y')=D^{-2}$. All
$\ell$ copies of both choice bits change by one, so
$\norm{y'-v^*}^2\ge2\ell$ and every valid growth constant satisfies
$g\le1/(2\ell D^2)$. The largest diagonal entry of $\nabla^2\Psi$ is ten,
attained by the digit-state variables. Therefore
\[
\kappa\ge20\ell D^2\ge80\ell c^2 .
\]
For $c=2^k$, $k\ge2$, there are $5k+7$ variables and coefficients bounded
by five, so $\kappa\ge80(k+1)4^k$ is exponential in the number of variables.
The instance has $O(k)$ monomials and bags, listed by variable indices of
$\Theta(\log k)$ bits, so $I=\Theta(k\log k)$ and
$\kappa\ge2^{\Omega(I/\log I)}$, consistent with
Corollary~\ref{lim:cor:nopolylog}.

Negative curvature is also large relative to growth. Let $\nu$ be the
smallest number with $\nabla^2\Psi\succeq-\nu I$. Encode the fractional
choice $(1/c,1)$; since $c\cdot c^{-1}+(c-1)\cdot1=c$, every square residual
can be kept at zero with all variables in $[0,1]$, and $\Psi(y)=\varrho/c$
with $\varrho=1-1/c$. The displacement $d=y-v^*$ annihilates every affine
residual, so $d^{\T}\nabla^2\Psi\,d=-2(\varrho^2+1)$, coming from the two
endpoint penalties. Every valid growth constant has
$g\le\varrho/(c\norm d^2)$, while $\nu\ge2(\varrho^2+1)/\norm d^2$. Hence
\[
\frac\nu g\ge2c\Bigl(\varrho+\frac1\varrho\Bigr)\ge4c .
\]
This bound on $\nu/g$ also holds for arbitrary positive weights on the
squares and a common positive weight on the two endpoint penalties: both
estimates scale with the penalty weight, and the square weights do not
enter. The reduction therefore does not contradict an algorithm
parameterized by $\nu/g$ either.
\end{remark}

\subsection{Proof of Proposition~\ref{prop:lbproduct}}\label{app:lbproduct}

\begin{lemma}[Isolation \cite{MulmuleyVaziraniVazirani1987}]\label{lem:isolation}
Let $\mathcal F$ be a nonempty family of subsets of a finite set $\mathcal U$,
and let $w:\mathcal U\to\{1,\dots,\rho\}$ be uniformly random. Then the minimum
of $w(Z)=\sum_{e\in Z}w(e)$ over $Z\in\mathcal F$ is attained by a unique set
with probability at least $1-\abs{\mathcal U}/\rho$.
\end{lemma}

\begin{proof}
For $e\in\mathcal U$ let $\alpha_e=\min\{w(Z):Z\in\mathcal F,e\notin Z\}
-\min\{w(Z\setminus\{e\}):Z\in\mathcal F,e\in Z\}$ (with
$\min\emptyset=+\infty$); $\alpha_e$ does not depend on $w(e)$. If two
distinct minimizing sets exist, some $e$ lies in exactly one of them, and
then $w(e)=\alpha_e$. For fixed $e$ this event has probability at most
$1/\rho$, because $w(e)$ is uniform and independent of $\alpha_e$. A union
bound over $e\in\mathcal U$ gives probability at most $\abs{\mathcal U}/\rho$
that the minimizing set is not unique.
\end{proof}

\begin{proof}[Proof of Proposition~\ref{prop:lbproduct}]
An instance of multicolored clique consists of $k\ge2$ color classes, each
identified with $\{1,\dots,N_0\}$, and, for each pair $c<d$ of classes, a
list $E_{cd}=((a^{cd}_1,b^{cd}_1),\dots,(a^{cd}_{m_{cd}},b^{cd}_{m_{cd}}))$
of distinct pairs, the edges between the classes; it asks for
$x_1,\dots,x_k$ with $(x_c,x_d)\in E_{cd}$ for all $c<d$. Let
$E=\bigsqcup_{c<d}E_{cd}$, choose independent uniform weights
$w(e)\in\{1,\dots,2\abs E\}$, and let $W_0=1+\binom k2\,2\abs E$. Use integer
variables $x_c\in[1,N_0]$ and, for each pair $c<d$ and $l=1,\dots,m_{cd}$,
binary $\zeta^{cd}_l$, $\sigma^{cd}_l$ and integer
$y^{cd}_l,{y'}^{cd}_l\in[0,N_0]$, with $\sigma^{cd}_0=y^{cd}_0={y'}^{cd}_0=0$.
Inside the brackets below the pair $c<d$ is fixed and the superscript $cd$
is omitted, so that $m=m_{cd}$, $a_l=a^{cd}_l$, $\zeta_l=\zeta^{cd}_l$, and so
on. Let
\begin{align*}
\Phi&=\sum_{c<d}\Bigl[\sum_{l}\bigl((\sigma_l-\sigma_{l-1}-\zeta_l)^2
+(y_l-y_{l-1}-a_l\zeta_l)^2+(y'_l-y'_{l-1}-b_l\zeta_l)^2\bigr)\\
&\qquad\qquad+(\sigma_m-1)^2+(y_m-x_c)^2+(y'_m-x_d)^2\Bigr],\qquad
\Psi_w=W_0\Phi+\sum_{c<d}\Bigl[\sum_lw(a_l,b_l)\zeta_l\Bigr] .
\end{align*}
\emph{Decomposition.} Take a central bag $\{x_1,\dots,x_k\}$; for each pair a
bag $\{x_c,x_d,\sigma_m,y_m,y'_m\}$ adjacent to it, followed by the path of
bags $\{\sigma_{l-1},y_{l-1},y'_{l-1},\zeta_l,\sigma_l,y_l,y'_l\}$,
$l=m,\dots,1$, with the constants $\sigma_0,y_0,y'_0$ omitted. Every term lies
in a bag, running intersection holds, and $p=\max\{k,7\}$.

\emph{Encoding.} $\Phi\ge0$ is an integer. $\Phi=0$ forces $\sigma$ to
increase from $0$ to $1$ in steps $\zeta_l\in\{0,1\}$, so exactly one
$\zeta_{l^*}=1$, and then $(x_c,x_d)=(y_m,y'_m)=(a_{l^*},b_{l^*})\in E_{cd}$.
Thus $\Phi=0$ exactly at encodings of multicolored cliques, and each clique has
exactly one encoding because the pairs in $E_{cd}$ are distinct. At an
encoding the linear term equals the weight of the clique's edge set, which is
at most $W_0-1$. Every point with $\Phi\ge1$ has $\Psi_w\ge W_0$. If a clique
exists and the minimum-weight clique is unique, $\Psi_w$ has a unique
minimizer, with value at most $W_0-1$.

\emph{Conditioning.} Suppose the minimum-weight clique is unique. The values
of $\Psi_w$ are integral and its minimizer $z^*$ is unique, so
$\Psi_w(z)-\OPT\ge1$ for $z\ne z^*$. All domains have length at most $N_0$,
and there are $n\le k+2k^2N_0^2$ variables, so
$\norm{z-z^*}^2\le nN_0^2$ and $g=1/(nN_0^2)$ is a valid growth constant.
The diagonal Hessian entries are at most $W_0(2+4N_0^2+2k)$. Hence
$\kappa\le W_0(2+4N_0^2+2k)nN_0^2$, and both $\kappa$ and $I$ are at most
$(kN_0)^{c_2}$ for an absolute constant $c_2$.

\emph{Reduction.} Run the assumed algorithm for
$f(p)((kN_0)^{2c_2})^{p/\psi(p)}$ steps. If it outputs a feasible $\hat z$
with $\Psi_w(\hat z)<W_0$, then $\Phi(\hat z)=0$ and $\hat z$ encodes a
clique; answer yes. Otherwise answer no. There are no false positives. If a
clique exists, Lemma~\ref{lem:isolation}, with $\mathcal U=E$, $\rho=2\abs E$
and $\mathcal F$ the family of edge sets of multicolored cliques (which
determine the cliques since $k\ge2$), makes the minimum-weight clique unique
with probability at least $1/2$, and then the algorithm halts in time with
$\Psi_w(\hat z)\le W_0-1/2$. For $k\ge7$ this is a one-sided-error randomized
algorithm for multicolored clique with running time
$f'(k)N_0^{2c_2k/\psi(k)+O(1)}$ for a computable $f'$.

Clique reduces to multicolored clique with the same $k$: take $k$ copies of
the vertex set as color classes, and join the copies of $u$ and $v$ in
different classes when $uv$ is an edge; the classes then have $N_0$ elements
for a graph on $N_0$ vertices. Under ETH, Clique on graphs with $N_0$ vertices
has no algorithm with running time $f(k)N_0^{o(k)}$ for a computable $f$
\cite{ChenEtAl2006}, \cite[Theorem~14.21]{CyganEtAl2015}; the exponent
$2c_2k/\psi(k)+O(1)$ is computable and $o(k)$. The proof of this lower bound
composes the sparsification lemma with reductions from 3-SAT to Clique, all
deterministic. Composing them and the reduction from Clique to multicolored
clique with the algorithm above, run twice to reduce the error probability
to $1/4$, gives a randomized algorithm for 3-SAT with error probability at
most $1/3$ and running time $2^{o(n')}$ on $n'$ variables, which contradicts
rETH.

The reduction uses only instances with $p=k$, and $k$ may be restricted to
$k\ge k_0$ for any fixed $k_0$, because finitely many values of $k$ do not
affect a bound $N_0^{o(k)}$. So the first claim holds even for algorithms whose
time bound is required only for $p\ge k_0$. For the last claim, increasing
$C$ if necessary, we may assume that $C$ is a positive integer. Suppose
$a(p)\le p/\psi(p)$ and put $\psi'(p)=\max\{1,\min\{\psi(p),p/C\}\}$. Then
$\psi'$ is computable, nondecreasing, unbounded and at least one, and for
$p\ge C$ it equals $\min\{\psi(p),p/C\}$, so that
$p/\psi'(p)\ge\max\{a(p),C\}$ and, because $\kappa,I\ge1$,
$f(p)\kappa^{a(p)}I^C\le f(p)(\kappa I)^{p/\psi'(p)}$. The first claim with
$k_0=\max\{7,C\}$ excludes this running time. Finally,
Theorem~\ref{thm:approx} with $q=1$ runs in time
$f(p,\kappa)(I+2)^5\le3^5f(p,\kappa)I^5$, where
$f(p,\kappa)=c_0(c_1p\sqrt\kappa)^p\kappa(1+\log_2\kappa)^2
\le4c_0(c_1p)^p\kappa^{p/2+2}$ because $(1+\log_2\kappa)^2\le4\kappa$. So
$a(p)=p/2+O(1)$ and $C=5$.
\end{proof}

\subsection{Proofs for Sections~\ref{sec:limits-messages}--\ref{sec:limits-constraints}}
\label{app:limits-proofs}

\begin{proof}[Proof of Proposition~\ref{lim:prop:messages}]
Part 1. $\xi_t$ lies in bags $t$ and $t+1$, and the endpoints $2^t-1$ have
$t$ bits.

Part 2. On the box $\xi_t\ge0$ and $z_t\ge0$, so
$\xi_{t-1}=(\xi_t-z_t-\rho_t)/2\le(\xi_t+\abs{\rho_t})/2$. By induction,
\[
0\le\xi_t\le2^{t-m}\xi_m+\sum_{j=t+1}^m2^{t-j}\abs{\rho_j} .
\]
The vector $(2^{t-m})_{t\le m}$ has squared norm at most $4/3$. The map
$y\mapsto(\sum_{j>t}2^{t-j}y_j)_t$ is the sum over $l\ge1$ of $2^{-l}$ times
the shift by $l$ places. Each shift has operator norm at most one, so the map
has operator norm at most $\sum_{l\ge1}2^{-l}=1$. Hence
$\norm\xi\le(2/\sqrt3)\xi_m+\norm\rho$ and
$\norm\xi^2\le\frac83\xi_m^2+2\norm\rho^2$. Also
$0\le z_t=\xi_t-2\xi_{t-1}-\rho_t\le\xi_t+\abs{\rho_t}$, so
$\norm z^2\le2\norm\xi^2+2\norm\rho^2$. Thus
$\norm\xi^2+\norm z^2\le3\norm\xi^2+2\norm\rho^2\le8\xi_m^2+8\norm\rho^2
\le8\Psi_m$.

Part 3. The second derivatives are $2+8=10$ for $\xi_t$, $t<m$; $2+2=4$ for
$\xi_m$; and $2-\frac14$ for $z_t$. Writing $\Psi_m$ as a sum of squares plus
$-\frac18\sum z_t^2$ plus linear terms gives
$\nabla^2\Psi_m\succeq-\frac14\sum_te_{z_t}e_{z_t}^{\T}\succeq-\frac14I$,
where $e_{z_t}$ is the unit vector of the coordinate $z_t$.

Part 4. Let $V_0(v)=V(v)-v^2
=\min\{\sum_t\rho_t^2+\frac18\sum_tz_t(1-z_t):\xi_m=v\}\ge0$. The minimum is
attained by compactness. Hence $V_0(v)=0$ exactly when some feasible point
has $\xi_m=v$, all $\rho_t=0$ and binary $z$. Then $\xi_t=2\xi_{t-1}+z_t$ and
$v=\sum_t2^{m-t}z_t\in\{0,\ldots,2^m-1\}$. Conversely, for an integer $j$ in
this range, its binary digits give $\xi_t=\lfloor j/2^{m-t}\rfloor\in
[0,2^t-1]$. So $V_0$ vanishes exactly at the $2^m$ integers of
$[0,2^m-1]$.

In the first representation, on an interval with nonempty interior,
$V_0$ coincides with a polynomial of degree at most two. It is not
identically zero, since $V_0>0$ off the integers. So it has at most two
zeros, and the interval contains at most two integers. A degenerate interval
contains at most one. Hence $2N_{\rm pc}\ge2^m$. In the second
representation, each integer $j$ has some $l$ with $\varphi_l(j)=V(j)$. Then
$\varphi_l-v^2\ge V_0\ge0$ on $[0,2^m-1]$ and vanishes at $j$. If
$\varphi_l-v^2\equiv0$, then $V_0\le0$, which is false; so $\varphi_l-v^2$
vanishes at most at two integers, and again $2N_{\rm pc}\ge2^m$.
\end{proof}

\begin{proof}[Proof of Proposition~\ref{lim:prop:setgrowth}]
For $x\in X$, $\dist(x,\mathcal S)^2\le\sum_i(x_i-x_1)^2$, and
$(x_i-x_1)^2\le(i-1)F(x)$ by Cauchy--Schwarz, which gives the set-growth
bound. For $n=2$, $\dist(x,\mathcal S)^2=(x_1-x_2)^2/2$. The coordinate second
derivatives are $2$ for $x_1,x_n$ and $4$ otherwise, so every valid $L_i$ is
at least $2$; with the smallest valid ones, $L=\max_iL_i\le4$, and $L=2$ if
$n=2$, so $\kappa_S\le2n(n-1)$.

(a) Fix $j$ and $v\in G_j$, and put $v_j=v$. Choose the remaining coordinates
outward along the path. If $v_k$ has been chosen and $i=k\pm1$ is the next
index, let $[\alpha,\alpha']$ be a grid interval of $G_i$ containing $v_k$, of
length $\delta_i$, and let $v_i$ be its endpoint nearest to $v_k$. Such an
interval exists because $G_i$ contains $0$ and $1$. Then
$(v_k-v_i)^2\le\delta_i^2/4$ and $w_i(v_i)\ge\delta_i$. Every edge of the
path joins some $i\ne j$ to the index chosen just before it, so
\[
Q(v)=F(v)-D(v)\le\sum_{i\ne j}\frac{\delta_i^2}4-\frac{L_j}8w_j(v)^2
-\sum_{i\ne j}\frac{L_i}8\delta_i^2\le-\frac{L_j}8w_j(v)^2,
\]
because $L_i\ge2$. Hence $m_j(v)\le Q(v)\le-\frac{L_j}8w_j(v)^2$, and an
endpoint $v$ of a longest grid interval gives the bound on $\beta$.

(b) Let $G$ be a grid of $X$, let $[a,a']$ be a grid interval of $G_i$, and
let $t\in[a,a']$. The point $t\mathbf1$ is optimal and lies in $X$, so
Proposition~\ref{prop:cellwise}(c) at $x=t\mathbf1$ gives
$\min\{m_i(a),m_i(a')\}\le F(t\mathbf1)=0\le U$, and the interval is
retained. By induction, every stage of such a run has the box $X$. Suppose that the
last stage certifies accuracy $\varepsilon$, that is, $U-\beta\le\varepsilon$
with $\beta$ the corrected minimum of the last grid and $U\ge\OPT=0$. By (a)
and $L_j\ge2$, $W_j^2/4\le-\beta\le U-\beta\le\varepsilon$, so
$W_j\le2\sqrt\varepsilon$. Covering $[0,1]$ by intervals of length at most
$2\sqrt\varepsilon$ needs at least $1/(2\sqrt\varepsilon)$ of them, so
$\abs{G_j}\ge1+1/(2\sqrt\varepsilon)$.

(c) Since $F(\hat x)\ge\OPT=0$, $\beta\ge-\varepsilon$. The grid $G^{(0)}$ is
a grid of $X^{(0)}=X$. Let $J=[a,a']$ be a grid interval of $G^{(0)}_i$ of
length $\delta>0$ that is not contained in $X^{(1)}_i$. Both $a$ and $a'$ are
endpoints of it, so $w_i(a),w_i(a')\ge\delta$, and (a) gives
$\min\{m^{(0)}_i(a),m^{(0)}_i(a')\}\le-\delta^2/4$. Condition (C1) of
Definition~\ref{def:cert} requires this minimum to be at least
$\beta\ge-\varepsilon$ (its second alternative needs $w(J)=0$, that is,
$\delta=0$), so $\delta\le2\sqrt\varepsilon$. If $k=0$, (C2) and (a) give
$-W_i^2/4\ge\min_{G^{(0)}}Q^{(0)}\ge\beta\ge-\varepsilon$ for every $i$.
\end{proof}

\begin{proof}[Proof of Proposition~\ref{prop:oraclebarrier}]
For $c\in X$ and $0<w\le1$ let
$f_{c,w}(x)=-\frac{w^2}6\bigl(1-\norm{x-c}^2/w^2\bigr)^3$ when
$\norm{x-c}\le w$ and $f_{c,w}(x)=0$ otherwise. In the variable $z=(x-c)/w$
with $\rho=\norm z^2$,
\[
\partial_if_{c,w}=w\,z_i(1-\rho)^2,\qquad
\partial_{ij}f_{c,w}=\delta_{ij}(1-\rho)^2-4z_iz_j(1-\rho).
\]
Value, gradient and Hessian vanish on $\rho=1$, so $f_{c,w}$ is $C^2$, and
$\partial_{ii}f_{c,w}=(1-\rho)(1-\rho-4z_i^2)\le1$. The minimizer $c$ is
unique with value $-w^2/6$. Inside the ball,
$f_{c,w}+\frac{w^2}6=\frac{w^2}6\bigl(1-(1-\rho)^3\bigr)\ge\frac{w^2}6\rho
=\frac{\norm{x-c}^2}6$, because $1-(1-\rho)^3-\rho=\rho(1-\rho)(2-\rho)\ge0$.
Outside it the gap is $w^2/6\ge\frac{w^2}{6n}\norm{x-c}^2$. Hence set growth
holds with $g_S=w^2/(6n)$. The zero function has $\argmin_Xf=X$ and satisfies
it with every $g_S$, in particular with $g_S=1$.

Run the algorithm on $f\equiv0$; let $Q$ be its number of queries and add
the output point, if any, to the query points. Suppose
$Q<\bigl(1/(2\sqrt{6\varepsilon})-1\bigr)^n-1$, and let
$M=\lfloor(Q+1)^{1/n}\rfloor+1$. Then $M^n>Q+1$ and $M<1/(2\sqrt{6\varepsilon})$.
Divide $X$ into $M^n$ closed subcubes of side $1/M$. The interior of one of
them contains none of the at most $Q+1$ points. Let $c$ be its center and
choose $w$ with $\sqrt{6\varepsilon}<w<1/(2M)$. The closed ball of radius $w$
lies in the open subcube, so every oracle answer for $f_{c,w}$ at the
recorded points equals the answer for $0$. By determinism the algorithm
produces the same transcript and output on $f_{c,w}$. An interval of width at
most $\varepsilon$ containing $0$ does not contain $-w^2/6<-\varepsilon$. An
output point outside the ball has gap $w^2/6>\varepsilon$ for $f_{c,w}$.
Either way the algorithm is wrong on an admissible function.
\end{proof}

\begin{proof}[Proof of Proposition~\ref{lim:prop:constraints}]
Summing the recurrences gives $\sum_{i=0}^ma_ix_i=a_0$. The choice $x_0=1$,
$x_i=0$ ($i\ge1$) is feasible. Since all $a_i>0$, it is the only feasible
choice with $x_0=1$. A choice with $x_0=0$ is feasible exactly when its
items sum to $a_0$: its partial sums then lie in $[0,1]$. For each binary $x$
the constraints determine $y$. Distinct feasible points therefore have
distinct binary vectors and distinct integer objective values. The dummy
value $2^{m+1}$ exceeds every value $\le2^m-1$ of a genuine subset. Hence
the minimizer is unique, and $x_0^*=0$ exactly for yes-instances. All
$2m+3$ coordinates lie in $[0,1]$, so $\norm{z-z^*}^2\le2m+3$, while
$F(z)-\OPT\ge1$ for $z\ne z^*$. An approximation within $1/2$ separates
$\OPT\le2^m-1$ from $\OPT=2^{m+1}$. The data have polynomial binary length,
and Subset Sum is NP-complete \cite{GareyJohnson1979}.
\end{proof}
